% !TeX root = thesis.tex

\chapter{Related Work}
\label{ch:related_work}

This chapter reviews prior work in five areas directly relevant to this
project: vision-language model benchmarks, multi-sensor datasets and temporal
alignment, day--night and low-light video understanding, LLM-assisted annotation
pipelines, and the vision-language models evaluated in this work.


\section{Vision-Language Model Benchmarks}
\label{sec:rw_vlm_benchmarks}

Evaluating Vision-Language Models (VLMs) has attracted substantial interest as
these models have advanced from image captioning~\cite{chen2015microsoft} to
complex visual reasoning. Early benchmarks such as VQA~\cite{antol2015vqa} and
GQA~\cite{hudson2019gqa} established visual question answering as a standard
evaluation paradigm. Subsequent work extended evaluation to compositional image
reasoning, including MMBench~\cite{liu2023mmbench} and
SEED-Bench~\cite{li2023seedbench}, and to video understanding through
ActivityNet-QA~\cite{yu2019activitynet}, NExT-QA~\cite{xiao2021nextqa},
Video-MME~\cite{fu2024videomme}, and MVBench~\cite{li2024mvbench}.

Video-based benchmarks evaluate abilities such as action recognition, temporal
ordering, object movement, state changes, and long-context reasoning. Egocentric
datasets such as EgoSchema~\cite{mangalam2023egoschema},
Ego4D~\cite{grauman2022ego4d}, and EgoNight~\cite{egonight2024} extend these
tasks to first-person recordings and human-object interactions. However, most
existing benchmarks rely primarily on conventional RGB images or videos and do
not systematically compare model behaviour across heterogeneous visual sensors.

Our benchmark provides four visual sensor modalities---RGB, infrared, event,
and depth. The sensor streams are
synchronized within each recording, while the daytime and nighttime recordings
represent semantically corresponding repetitions of the same activity. This
design enables controlled comparison of VLM robustness across sensor modality,
illumination condition, and input representation.


\section{Multi-Sensor Datasets and Temporal Alignment}
\label{sec:rw_multisensor}

Non-standard visual sensors provide physical properties that complement
traditional RGB cameras. Event cameras capture asynchronous per-pixel brightness
changes at very high temporal resolution, making them useful for rapid motion
and low-light scenes. Canonical event-camera datasets include
N-Caltech101~\cite{orchard2015converting} and the DAVIS Event
Dataset~\cite{mueggler2017event}. Infrared datasets such as
FLIR~\cite{flir2019adas} and KAIST~\cite{hwang2015multispectral} have
demonstrated the value of non-visible-spectrum sensing under difficult
illumination and weather conditions.

Depth data contributes information about distance, near--far ordering,
occlusion, free space, and three-dimensional scene structure. In addition to
sensor-recorded depth, monocular depth-estimation methods such as
Marigold~\cite{ke2024repurposing} can generate dense depth maps from RGB
images.

Temporal alignment is required before these heterogeneous streams can form
shared evidence units. Dynamic Time Warping
(DTW)~\cite{berndt1994using}, cross-correlation, and optical flow have been
used to align sensor streams in robotics, audio-visual learning, and action
recognition~\cite{arandjelovic2017look}. Our pipeline combines these techniques
according to the signal characteristics: cross-correlation estimates fixed
offsets, optical-flow activity provides a common motion representation, and DTW
models non-linear timing differences before task-level slicing is applied.


\section{Day--Night and Low-Light Video Understanding}
\label{sec:rw_day_night}

Changes in illumination introduce a substantial domain shift for visual
recognition systems. Nighttime RGB recordings may contain reduced contrast,
weaker colour information, motion blur, sensor noise, and poorly defined object
boundaries. Although datasets such as EgoNight~\cite{egonight2024} address
first-person understanding under low illumination, relatively few benchmarks
evaluate open-ended video question answering using corresponding daytime and
nighttime recordings of the same activity.

Nighttime observations should not always be interpreted only as degraded
versions of daytime scenes. Artificial lights, illuminated displays,
reflections, and bright regions may become more salient at night. In addition,
event cameras can preserve motion cues under difficult lighting, while infrared
sensors may retain object and scene structure when visible-spectrum RGB
information becomes unreliable.

Our benchmark therefore evaluates four contrastive outcomes: daytime correct
and nighttime incorrect, nighttime correct and daytime incorrect, both correct,
and both incorrect. These categories distinguish illumination-specific failures
from generally difficult questions. They also reveal whether a particular
sensor provides stronger evidence under specific lighting conditions rather
than assuming that daytime input is always superior.


\section{Automated Dataset Annotation with LLMs}
\label{sec:rw_auto_annotation}

The availability of instruction-tuned language and vision models has introduced
a scalable approach to dataset construction, reducing dependence on expensive
human crowd-sourcing. LLaVA~\cite{liu2023llava} demonstrated that GPT-4-generated
instruction data could support high-quality visual instruction tuning.
DataComp~\cite{gadre2023datacomp} showed that dataset curation strategy is as
important as raw dataset size, while LAION-COCO~\cite{schuhmann2022laion} and
ShareGPT4V~\cite{chen2023sharegpt4v} illustrated large-scale caption and
instruction generation.

Multi-sensor annotation requires modality-specific generation strategies because
different sensors represent different forms of evidence. RGB questions may
focus on appearance and object identity, event questions on motion and temporal
change, depth questions on geometry and spatial relations, and infrared questions
on low-light visibility. A single generic prompt is therefore insufficient for
constructing a modality-aware benchmark.

Automatically generated annotations may contain hallucinations, vague or
unanswerable questions, unsupported answers, and contradictions with the
available evidence. Distillation-based filtering~\cite{west2022symbolic} and
LLM judges~\cite{zheng2023judging} have been used to improve generated-data
quality. Our pipeline applies a lightweight semantic verifier based on Gemini
3.1 Flash Lite, which evaluates each QA candidate using five binary quality
criteria and a hallucination-risk score. Only candidates that pass all validity
checks are retained before multimodal evidence fusion.


\section{Vision-Language Models Evaluated in This Work}
\label{sec:rw_vlms}

To demonstrate the utility of the constructed multi-sensor benchmark, we
evaluate three recent open-weight VLM families at approximately 4B and 8B
parameter scales. Qwen3-VL~\cite{wang2024qwen2vl} supports dynamic-resolution
visual processing and detailed spatial and temporal reasoning.
InternVL3.5~\cite{chen2024internvl} uses a high-resolution InternViT encoder and
a progressive vision-language alignment strategy, while
Molmo2~\cite{deitke2024molmo} combines a vision encoder with an efficient
decoder-only language model and video-processing components.

The models are evaluated using fixed-frame inputs and, where supported by the
corresponding implementation, native-video inputs. Fixed-frame evaluation
provides a controlled and computationally predictable representation, whereas
native-video input allows models to use their internal temporal-processing
mechanisms. This comparison tests whether native-video processing consistently
offers an advantage over representative frame sampling.

Evaluating both 4B and 8B variants also allows us to investigate whether a
larger parameter count produces consistent improvements across RGB, infrared,
event, and depth inputs. The experiments therefore examine model robustness
along four dimensions: sensor modality, illumination condition, model scale,
and evidence-presentation format.